\subsection*{向量}
\begin{enumerate}
    \item 已知 $\bigtriangleup {OAB}$ 中, $\vv{OA} = \vv{a},\vv{OB} = \vv{b}$ ,满足 $\left| \vv{a}\right|  = \left| \vv{b}\right|  = \left| {\vv{a} - \vv{b}}\right|  = 3$ ,求 $\left| {\vv{a} + \vv{b}}\right|$ 与 $\bigtriangleup {OAB}$ 面积.
    \begin{jieda}
        由已知得 $\left| \vv{OA}\right|  = \left| \vv{OB}\right|$ ,以 ${OA},{OB}$ 为邻边作平行四边形 ${OACB}$ ,则可知其为菱形, 且 $\vv{OC} = \vv{a} + \vv{b},\vv{BA} = \vv{a} - \vv{b}$ ,由于 $\left| \vv{a}\right|  = \left| \vv{b}\right|  = \left| {\vv{a} - \vv{b}}\right|  = 3$ ,则 $\left| {OA}\right|  = \left| {OB}\right|  = \left| {BA}\right|  = 3$，根据几何关系可知$\left| {\vv{a} + \vv{b}}\right|  = \left| {OC}\right|  = 3\sqrt{3}$，${S}_{\bigtriangleup {OAB}} = \dfrac{1}{2}\left| {OA}\right|  \cdot  \left| {OB}\right| \sin \dfrac{\pi }{3} = \dfrac{\sqrt{3}}{4} \times  {3}^{2} = \dfrac{9\sqrt{3}}{4}$
        
    
        \begin{center}
        \includegraphics[width=0.3\textwidth]{bodyxb/src/bo_d5dqkaref24c73bjk040_0_240_1570_372_250_0.jpg}
        \end{center}

    \end{jieda}
    \item 设平面上有四个互异的点$A,B,C,D$，且$(\vv{DB} + \vv{DC} -2\vv{DA}) \cdot (\vv{AB} - \vv{AC}) = 0$，则$\triangle ABC$的形状是\dotfill{(\kh{B})}
    \xx{直角三角形}{等腰三角形}{等腰直角三角形}{等边三角形}
    \begin{jieda}
        $(\vv{DB} + \vv{DC} -2\vv{DA}) \cdot (\vv{AB} - \vv{AC}) = 0$即$(\vv{AB} + \vv{AC}) \cdot \vv{CB} = 0$，即$BC$边中线与$BC$垂直，根据几何关系可知$AB = AC$，即$\triangle ABC$为等腰三角形.
    \end{jieda}
    \item 已知$A,O,B$是不共线的三点，且$C,D$满足$\vv{OC} = 2 \vv{OA}, \vv{OD} = 3 \vv{OB}$，$AD$与$BC$交于点$M$，设$\vv{OM} = x \vv{OA} + y \vv{OB}$，试求$x, y$.
    \begin{jieda}
        $\vv{OM} = x \vv{OA} + y \vv{OB} \Rightarrow \left \{ \begin{array}{l}
                \vv{OM} = \dfrac{x}{2}\vv{OC} + y\vv{OB}  \\
                \vv{OM} = x \vv{OA} + \dfrac{y}{3}\vv{OD} 
        \end{array} \right.$，因为$A, M, D$共线，$C, M, B$共线，根据等和线结论有$\left \{ \begin{array}{l}
                \dfrac{x}{2} + y = 1  \\
                x + \dfrac{y}{3} = 1 
        \end{array} \right.$，解得$x = \dfrac{4}{5}, y = \dfrac{3}{5}$.
    \end{jieda}

    \item 若向量$\vv{a} = \left(\dfrac{3}{2}, sin\alpha\right), \vv{b} = \left(cos\alpha, \dfrac{1}{3}\right)$，且$\vv{a} // \vv{b}$，则锐角$\alpha = $\blkda{$\dfrac{\pi}{4}$}
    \begin{jieda}
        根据平面向量平行的充要条件的坐标表示有$sin\alpha \cdot cos\alpha = \dfrac{1}{2}$，故$sin2\alpha = 1$，因此$\alpha = \dfrac{\pi}{4}$
    \end{jieda}

    \item 设$\vv{a}, \vv{b}, \vv{c}$是任意的非零向量，且两两不共线，则下列命题中正确的是\blkda{\ding{173} \ding{175}}:
    \begin{dingautolist}{172}
        \item $(\vv{a} \cdot \vv{b}) \cdot \vv{c} - (\vv{c} \cdot \vv{a}) \cdot \vv{b} = \vv{0}$
        \item $|\vv{a}| - |\vv{b}| < |\vv{a} - \vv{b}|$
        \item $(\vv{b} \cdot \vv{c}) \cdot \vv{a} - (\vv{a} \cdot \vv{c}) \cdot \vv{b}$不与$\vv{c}$垂直
        \item $(3\vv{a} + 2\vv{b}) \cdot (3\vv{a} - 2\vv{b}) = 9|\vv{a}|^2 - 4|\vv{b}|^2$
    \end{dingautolist}


    \item 已知向量 $\vv{AB} = \left( {k,1}\right) ,\vv{AC} = \left( {1,0}\right)$ ， $\bigtriangleup  {ABC}$ 是直角三角形，则 $k =$ \blkda{0 或 1}
    \begin{jieda}
        因为 $\vv{AB} = \left( {k,1}\right) ,\vv{AC} = \left( {1,0}\right)$ ,所以 $\vv{CB} = \vv{AB} - \vv{AC} = \left( {k - 1,1}\right)$ ,
        
        若 $\vv{AB} \bot  \vv{AC}$ ,则 $\vv{AB} \cdot  \vv{AC} = k = 0$ ,所以 $k = 0$ ,符合题意;
        
        若 $\vv{AB} \bot  \vv{CB}$ ,则 $\vv{AB} \cdot  \vv{CB} = k\left( {k - 1}\right)  + 1 = 0$ ,此方程无解;
        
        若 $\vv{AC} \bot  \vv{CB}$ ,则 $\vv{AC} \cdot  \vv{CB} = k - 1 = 0$ ,所以 $k = 1$ ,符合题意;
        
        综上, $k = 0$ 或 1 .
    \end{jieda}

    \item 如图所示,已知 $\bigtriangleup  {ABC}$ 中，点 $P,Q,R$ 依次是边 ${BC}$ 上的三个四等分点，若 ${BC} = 8$ ， $\vv{AP} \cdot  \vv{AR} = {20}$ ，则 $\vv{AB} \cdot  \vv{AC} =$ \blkda{8}.
    
    \begin{center}
    \includegraphics[width=0.2\textwidth]{bodyxb/src/3_241_1345_346_180_0.jpg}
    \end{center}
    \begin{jieda}
        $20 = \vv{AP} \cdot  \vv{AR} = |\vv{AQ}|^2 - |\vv{PQ}|^2 = |\vv{AQ}|^2 - 4$

        $\vv{AB} \cdot  \vv{AC} = |\vv{AQ}|^2 - |\vv{BQ}|^2 = |\vv{AQ}|^2 - 16 = 8$
    \end{jieda}
\end{enumerate}